\documentclass[11pt,a4paper,fontset=windows]{ctexart}
% ===== 公共导言区：高等代数（一）期中模拟卷 =====
\usepackage{amsmath,amssymb}
\usepackage{unicode-math}
\setmathfont{TeX Gyre Termes Math}
\usepackage[margin=2.05cm,top=2.35cm,bottom=2.0cm,headsep=11pt,includefoot]{geometry}
\usepackage{fancyhdr}
\usepackage{lastpage}
\usepackage{enumitem}
\usepackage{array}
\usepackage{bm}
\usepackage{xcolor}
\usepackage{zhnumber}
\usepackage{needspace}

\definecolor{hdgray}{gray}{0.32}
\definecolor{lngray}{gray}{0.5}
\definecolor{notebg}{gray}{0.955}

\pagestyle{fancy}
\fancyhf{}
\renewcommand{\headrulewidth}{0.4pt}
\renewcommand{\footrulewidth}{0pt}
\fancyhead[L]{\small\color{hdgray}\songti 高等代数（一）\ 期中模拟试卷（\thepp）}
\fancyhead[R]{\small\color{hdgray}\songti 考试时间：120 分钟\quad 满分：100 分}
\fancyfoot[C]{\small\color{lngray}\songti 第 \thepage\ 页\quad 共 \pageref{LastPage}\ 页}

\setlength{\parindent}{0pt}
\setlength{\parskip}{2pt}
\linespread{1.16}

% 题目列表 1—8 连续编号
\newlist{qs}{enumerate}{1}
\setlist[qs]{label=\textbf{\arabic*.},ref=\arabic*,leftmargin=2.55em,
  labelsep=0.55em,itemsep=23pt,topsep=6pt,parsep=0pt}
\newlist{note}{enumerate}{1}
\setlist[note]{label=(\arabic*),leftmargin=1.75em,itemsep=1.5pt,topsep=2.5pt,parsep=0pt}
\newlist{sub}{enumerate}{1}
\setlist[sub]{label=(\arabic*),leftmargin=1.9em,itemsep=3pt,topsep=3pt,parsep=0pt}

\newcommand{\secA}[1]{\Needspace{9\baselineskip}\par\vspace{10pt}{\large\bfseries\songti #1}\par\vspace{3pt}\nopagebreak}
\newcommand{\blank}[1][2.6cm]{\underline{\hspace{#1}}}
\newcounter{pp}
\renewcommand{\thepp}{\chinese{pp}}

\setcounter{pp}{1}
\begin{document}

\begin{center}
{\LARGE\bfseries\songti 高等代数（一）期中模拟试卷（一）}\\[7pt]
{\normalsize\songti 适用范围：第一章\ 多项式（\S1\,--\,\S9）；第二章\ 行列式（\S2.1\,--\,\S2.7）}\\[3pt]
{\small\color{hdgray}\songti 多元多项式、对称多项式不在考查范围}
\end{center}
\vspace{-4pt}

\begin{center}
\begin{tabular}{@{}p{0.335\textwidth}p{0.335\textwidth}p{0.26\textwidth}@{}}
\songti 姓名：\blank[2.5cm] & \songti 学号：\blank[2.5cm] & \songti 得分：\blank[1.9cm]
\end{tabular}
\end{center}

\vspace{2pt}
\hrule height 0.5pt
\begin{center}
\setlength{\fboxsep}{6pt}%
\colorbox{notebg}{\parbox{0.94\textwidth}{\small\songti
\textbf{考生须知：}
\begin{note}
\item 本试卷共 8 题，满分 100 分，考试时间 120 分钟。
\item 计算题须写出必要的计算过程，证明题须给出完整的推理步骤；只写答案不给分。
\item 所有答案一律写在答题纸上，写在试卷上无效。
\end{note}}}
\end{center}
\hrule height 0.5pt
\vspace{4pt}

\secA{一、计算题（每小题 10 分，共 40 分）}
\begin{qs}
\item 设
\[
f(x)=x^{4}+x^{3}-x-1,\qquad g(x)=x^{3}-x^{2}+x-1 .
\]
求 $(f(x),\,g(x))$，并求出多项式 $u(x),\,v(x)\in\mathbb{Q}[x]$，使得
\[
u(x)f(x)+v(x)g(x)=(f(x),\,g(x)).
\]

\item 计算 $n$ 阶行列式
\[
D_n=\begin{vmatrix}
13 & 6 & 0 & \cdots & 0\\
6 & 13 & 6 & \cdots & 0\\
0 & 6 & 13 & \ddots & \vdots\\
\vdots & \ddots & \ddots & \ddots & 6\\
0 & \cdots & 0 & 6 & 13
\end{vmatrix}.
\]

\item 设 $a_1a_2\cdots a_n\neq 0$。计算 $n$ 阶行列式
\[
D_n=\begin{vmatrix}
1+a_1 & 1 & 1 & \cdots & 1\\
1 & 1+a_2 & 1 & \cdots & 1\\
1 & 1 & 1+a_3 & \cdots & 1\\
\vdots & \vdots & \vdots & \ddots & \vdots\\
1 & 1 & 1 & \cdots & 1+a_n
\end{vmatrix}.
\]

\item 求多项式
\[
f(x)=x^{4}-x^{3}-7x^{2}+x+6
\]
在 $\mathbb{Q}[x]$ 中的标准分解式，并写出 $f(x)$ 在 $\mathbb{R}[x]$ 与 $\mathbb{C}[x]$ 中的分解式。
\end{qs}

\secA{二、证明题（每小题 15 分，共 60 分）}
\begin{qs}[resume*]
\item 记
\[
F=\{a+b\sqrt{2}\mid a,b\in\mathbb{Q}\},\qquad R=\{a+b\sqrt{2}\mid a,b\in\mathbb{Z}\}.
\]
\begin{sub}
\item 证明 $F$ 是一个数域；
\item 证明 $R$ 不是数域。
\end{sub}

\item 设 $m,n$ 为正整数，$d=(m,\,n)$ 为 $m$ 与 $n$ 的最大公因数。证明：
\[
\bigl(x^{m}-1,\;x^{n}-1\bigr)=x^{d}-1 .
\]

\item 设 $a_1,a_2,\ldots,a_n$ 是数域 $P$ 中 $n$ 个互不相同的数，令
\[
\ell_i(x)=\prod_{j\neq i}\frac{x-a_j}{a_i-a_j}\qquad (i=1,2,\ldots,n).
\]
\begin{sub}
\item 证明 $\displaystyle\sum_{i=1}^{n}\ell_i(x)\equiv 1$；
\item 证明：若 $g(x)\in P[x]$ 且 $\deg g<n$，则 $\displaystyle\sum_{i=1}^{n}g(a_i)\,\ell_i(x)=g(x)$。
\end{sub}

\item 设 $n\ge 2$，$n$ 阶行列式
\[
D_n=\begin{vmatrix}
x & a & a & \cdots & a\\
a & x & a & \cdots & a\\
a & a & x & \cdots & a\\
\vdots & \vdots & \vdots & \ddots & \vdots\\
a & a & a & \cdots & x
\end{vmatrix},\qquad a\neq x .
\]
\begin{sub}
\item 证明 $D_n=(x-a)^{n-1}\bigl(x+(n-1)a\bigr)$；
\item 求方程 $D_n=0$ 在 $\mathbb{C}$ 中的全部根。
\end{sub}
\end{qs}

\end{document}
